% Get better typography
% \usepackage[protrusion=true,expansion=true]{microtype}		

% Small margins
%\usepackage[top=2cm, bottom=2cm, left = 1cm, right = 1cm,columnsep=2

% For basic math, align, fonts, etc.
\usepackage{amsmath}
% \usepackage{unicode-math}
\usepackage{amsthm}
\usepackage{amssymb}
\usepackage{bbm}
\usepackage{mathtools}
\usepackage{mathrsfs}
\mathtoolsset{showonlyrefs}

\newtheorem{defn}{Definition}[]
\newtheorem{thm}{Theorem}[]
\newtheorem{claim}{Claim}[]
\newtheorem{lemma}{Lemma}[]
\newtheorem{prop}{Property}[]
\newtheorem{ass}{Assumption}[]
\newtheorem{cor}{Corollary}[]
\newtheorem{rem}{Remark}[]

% For creating sub-groups of Assumptions/Theorems etc.
\newtheorem{assumpA}{Assumption}
\renewcommand\theassumpA{A\arabic{assumpA}}

\usepackage{courier} % For \texttt{foo} to put foo in Courier (for code / variables)

\usepackage{lipsum} % For dummy text

% For images
\usepackage{graphicx}
% \usepackage{subcaption}
\usepackage[space]{grffile} % For spaces in image names

% For bibliography
% \usepackage[round]{natbib}

% For links (e.g., clicking a reference takes you to the phy)
\usepackage{hyperref}

\usepackage{theoremref}
% For color
\usepackage{graphicx}
\usepackage{wrapfig}
\usepackage{xcolor}
\definecolor{dark-red}{rgb}{0.4,0.15,0.15}
\definecolor{dark-blue}{rgb}{0,0,0.7}
\hypersetup{
    colorlinks, linkcolor={dark-blue},
    citecolor={dark-blue}, urlcolor={dark-blue}
}

% Macros
\newcommand{\emma}[1]{\textcolor{purple}{[Emma: #1]}}

\usepackage[colorinlistoftodos]{todonotes}

% Macros
\newcommand{\yash}[1]{\todo[color=blue!40]{Yash: #1}}

\newcommand{\red}[1]{\textcolor{red}{[#1]}}
\newcommand{\blue}[1]{\textcolor{blue}{[#1]}}





% \usepackage{algpseudocode}
% \usepackage[noend,linesnumbered,ruled]{algorithm2e}
% \usepackage[vlined,linesnumbered,ruled,algo2e]{algorithm2e}
% \usepackage[algo2e]{algorithm2e} 
% \SetKwProg{Fn}{Function}{}{}
\let\oldnl\nl% Store \nl in \oldnl
\newcommand{\nonl}{\renewcommand{\nl}{\let\nl\oldnl}}% Remove line number for one line


% Miscelaneous headers
\usepackage{multicol}
% \usepackage{enumitem}
\usepackage[utf8]{inputenc} % allow utf-8 input
\usepackage[T1]{fontenc}    % use 8-bit T1 fonts
\usepackage{hyperref}       % hyperlinks
\usepackage{url}            % simple URL typesetting
\usepackage{booktabs}       % professional-quality tables
\usepackage{nicefrac}       % compact symbols for 1/2, etc.
\usepackage[parfill]{parskip}

% \DeclareMathOperator*{\argmin}{arg\,min}
% \DeclareMathOperator*{\argmax}{arg\,max}
\DeclareMathOperator*{\ind}{\perp\!\!\!\perp}
\DeclareMathOperator*{\notind}{\not\!\perp\!\!\!\perp}
\DeclareMathOperator*{\cov}{\texttt{Cov}}
% \DeclareMathOperator*{\E}{\mathbb E}
% \DeclareMathOperator*{\R}{\mathbb R}


% Definitions of handy macros can go here
\newcommand{\dataset}{{\cal D}}
\newcommand{\fracpartial}[2]{\frac{\partial #1}{\partial  #2}}
\newcommand{\pibopt}{\pi_{b^\star}}
\newcommand{\pib}{{\pi_b}}
\newcommand{\pitheta}{{\pi_\btheta}}
\newcommand{\deltheta}{{\frac{\partial}{\partial \btheta}}}
\newcommand{\delthetatwo}{{\frac{\partial^2}{\partial^2 \btheta}}}
\newcommand{\pieval}{{\pi_e}}
\newcommand{\btheta}{{\boldsymbol\theta}}
% \newcommand{\data}{{\mathcal{D}}}
\newcommand{\ope}{{\operatorname{OPE}}}
\newcommand{\pe}{{\operatorname{PE}}}
\newcommand{\variance}{{\operatorname{Var}}}
\newcommand{\bias}{{\operatorname{Bias}}}
\newcommand{\mse}{{\operatorname{MSE}}}
\newcommand{\loss}{{\mathcal{L}}}
\newcommand{\horizon}{{l}}

% Auto sized brackets
\DeclarePairedDelimiter\autobracket{(}{)}
\newcommand{\br}[1]{\autobracket*{#1}}

\DeclarePairedDelimiter\autobrackett{[}{]}
\newcommand{\brr}[1]{\autobrackett*{#1}}

\DeclarePairedDelimiter\autobrackettt{\{}{\}}
\newcommand{\brrr}[1]{\autobrackettt*{#1}}
